\esSezione{Sistema PO/PC su board}\label{sec:esercizio6-2}

\begin{traccia}
Sintetizzare ed implementare su board il sistema PO/PC dell'esercizio precedente, rendendone osservabile il funzionamento: la scansione della ROM deve avanzare un passo alla volta a ogni pressione di un bottone, e il contenuto della MEM, insieme all'indirizzo corrente e al dato prodotto dalla macchina $M$, deve poter essere letto sui LED.
\end{traccia}

\progetto

Il sistema è lo stesso presentato nell'Esercizio~\ref{sec:esercizio6-1}: parte operativa composta da \texttt{counter\_mod}, \texttt{rom}, \texttt{onecount} e \texttt{mem}, unità di controllo \texttt{control\_unit}. Alla velocità del clock a \SI{100}{MHz} della board l'intera scansione terminerebbe in pochi cicli e non sarebbe osservabile; per questo si è scelto di sincronizzare l'avanzamento sulla pressione di un bottone, anziché rallentare il clock. Si aggiungono quindi due elementi: un ingresso \texttt{step} nell'unità di controllo, che consente il passaggio alla lettura successiva solo quando vale '1', e un top module \texttt{sistema\_tl} che istanzia il \texttt{debouncer} (già presentato nell'Esercizio~\ref{sec:esercizio3-2}, lst.~\ref{lst:esercizio3-2:debouncer}) per trasformare ogni pressione di \texttt{BTNC} in un impulso di un solo ciclo.

Per rendere visibile lo stato interno, il componente \texttt{sistema} espone anche due uscite di debug, \texttt{addr\_dbg} e \texttt{m\_dbg}, collegate a \texttt{addr} e \texttt{m\_data}. La \texttt{rom} (lst.~\ref{lst:esercizio6-1:rom}), la macchina $M$ \texttt{onecount} (lst.~\ref{lst:esercizio6-1:onecount}), la memoria \texttt{mem} (lst.~\ref{lst:esercizio6-1:mem}) e il contatore \texttt{counter\_mod} (lst.~\ref{lst:esercizio5-1:counter-mod}) restano invariati e non vengono ripresentati.

\implementazione

Presentato il criterio di avanzamento, vediamo le modifiche all'unità di controllo rispetto al listato~\ref{lst:esercizio6-1:cu}. Alla lista delle porte si aggiunge \texttt{step} (ingresso, 1 bit), e la stessa \texttt{step} entra nella sensitivity list del process combinatorio \texttt{f\_controllo}. Lo stato \texttt{S\_READ} ora condiziona tutto a \texttt{step}: finché vale '0' il comando \texttt{rd} resta a '0' (valore di default) e la macchina rimane nello stato; quando vale '1', \texttt{rd} è portato a '1' e lo stato prossimo diventa \texttt{S\_WRITE}. Gli altri stati e il process sincrono \texttt{mem\_stato}, con reset sincrono verso \texttt{S\_IDLE}, non cambiano. Poiché l'impulso del debouncer dura un solo ciclo, ogni pressione produce esattamente una lettura della ROM e una scrittura della MEM.

\vhdlfile[firstline=23,lastline=60]{esercizio6-2/control_unit.vhd}{Process combinatorio dell'unità di controllo con \texttt{step}}{lst:esercizio6-2:control-unit}

Nel componente \texttt{sistema} (listato~\ref{lst:esercizio6-1:sistema}) la porta \texttt{step} è inoltrata all'unità di controllo nel port mapping dell'istanza \texttt{u\_cu}, e sono aggiunte le uscite \texttt{addr\_dbg} (\texttt{A} bit) e \texttt{m\_dbg} (4 bit), assegnate con due istruzioni concorrenti a \texttt{addr} e \texttt{m\_data}. L'indirizzo mostrato è quello corrente del contatore, che avanza dopo la scrittura: nel commento del codice è indicato come indirizzo successivo.

\vhdlfile[firstline=10,lastline=22]{esercizio6-2/sistema.vhd}{Entità \texttt{sistema} con \texttt{step} e uscite di debug}{lst:esercizio6-2:sistema}

\vhdlfile[firstline=52,lastline=57]{esercizio6-2/sistema.vhd}{Istanza di \texttt{control\_unit} e uscite di debug}{lst:esercizio6-2:sistema-dbg}

Il top module \texttt{sistema\_tl} ha come porte \texttt{CLK100MHZ}, lo switch \texttt{SW} (16 bit), i bottoni \texttt{BTNC} e \texttt{BTND} e i LED \texttt{LED} (16 bit). L'architettura \texttt{structural} assegna il reset a \texttt{BTND} (attivo alto), istanzia il \texttt{debouncer} con \texttt{BTNC} in ingresso e l'impulso \texttt{step\_p} in uscita, e istanzia \texttt{sistema} con \texttt{N} pari a 16 e \texttt{a} pari a 4. Il mapping degli ingressi e delle uscite è il seguente:
\begin{enumerate}
  \item \texttt{SW(15)}: segnale \texttt{start};
  \item \texttt{SW(3 downto 0)}: indirizzo di lettura \texttt{raddr} della MEM;
  \item \texttt{BTNC}: avanzamento di un passo (\texttt{step}), filtrato dal debouncer;
  \item \texttt{BTND}: reset;
  \item \texttt{LED(3 downto 0)}: dato letto dalla MEM (\texttt{dout});
  \item \texttt{LED(7)}: segnale \texttt{done};
  \item \texttt{LED(11 downto 8)}: dato prodotto dalla macchina $M$ (\texttt{m\_dbg});
  \item \texttt{LED(15 downto 12)}: indirizzo del contatore (\texttt{addr\_dbg}).
\end{enumerate}
I LED da 4 a 6 sono forzati a \texttt{"000"}.

\vhdlfile{esercizio6-2/sistema_tl.vhd}{Implementazione del top module \texttt{sistema\_tl}}{lst:esercizio6-2:sistema-tl}

\simulazione

% TODO: testbench e simulazione assenti nel progetto Vivado dell'esercizio 6.2; la verifica del sistema senza step è nell'Esercizio 6.1.
Per questo esercizio non è stato sviluppato un testbench dedicato: la verifica funzionale della parte operativa e della parte di controllo è quella dell'Esercizio~\ref{sec:esercizio6-1}; la versione con 	exttt{step} si differenzia solo per il criterio di avanzamento.

\sintesiboard

Sulla board Nexys A7-100T, dopo il reset con \texttt{BTND}, si porta \texttt{SW(15)} a '1' per avviare il sistema; ogni pressione di \texttt{BTNC} produce un passo di scansione, durante il quale \texttt{LED(11 downto 8)} riporta il numero di '1' del byte letto dalla ROM e \texttt{LED(15 downto 12)} l'indirizzo del contatore. Terminate le sedici locazioni, \texttt{LED(7)} si accende; a quel punto il contenuto della MEM si consulta selezionando l'indirizzo con \texttt{SW(3 downto 0)} e leggendo il valore su \texttt{LED(3 downto 0)}.
% TODO: foto della scheda non disponibili nel progetto; nessun esito su board riportato.

Nel file dei constraint sono state decommentate le seguenti righe.

\xdcfile{esercizio6-2/vincoli.xdc}{Constraint}{lst:esercizio6-2:vincoli}
